%! Author = Omar Iskandarani
%! Companion to Swirl-String-Theory Canon-v0.8.1 (included from \texttt{SST\_CANON-v0.8.1.tex} before the manual bibliography).

\section{Research-track sectors migrated from the main canon}

\noindent
This companion document collects those sectors that remain useful for continuity and future development but are too long, too provisional, or too phenomenology-specific for the core body of \textit{Swirl-String-Theory\_Canon-v0.8.1}. Their migration out of the main canon is editorial rather than negative: each subsection remains part of the broader SST research program, but none is required for the minimal logical closure of the core canon.

\medskip
\noindent
\textbf{Editorial note (v0.8.1):} the maximal swirl tension sector ($F_{\rm swirl}^{\max}$, Rydberg bridge identities, Gibbons-class maximum-tension comparison, and the threefold interpretation table) is developed in the core document under \verb|SST_CANON-v0.8.1.tex|, subsection \texttt{Maximal Swirl Tension} (label \texttt{sec:Fmax}); it is not duplicated here.

\medskip

\subsection{Thermodynamic Radiation Interface}
A recurring legacy idea is that a local swirl-energy density may define an effective temperature field and thereby an emergent radiation sector. Let
\begin{align}
    T_{\mathrm{eff}} = T_{\mathrm{eff}}\big(\rhoE,\, \omega,\, \SwirlClock\big)
\end{align}
denote an unspecified constitutive relation. A minimal legacy interface then asks whether Planck-type spectra can emerge from a medium whose local excitation scale is set by $T_{\mathrm{eff}}$.

\textbf{[ORTHODOX]}: blackbody spectra are governed by Planck's law once a valid temperature variable is defined.

\textbf{[SPECULATIVE]}: SST does not yet derive the constitutive relation above, so this sector remains a research program rather than canonically closed physics.

\subsection{Minimal Coupling Interface to QED}
Early research notes repeatedly explored whether an effective vector potential could be built from circulation variables. The most conservative placeholder takes the form
\begin{align}
    \nabla \rightarrow \nabla - i\frac{m}{\hbar}A_{\mathrm{eff}},
    \qquad
    A_{\mathrm{eff}} = \chi \, v,
\end{align}
where $\chi$ is a scalar weight and $v$ is the local carrier velocity field.

\textbf{[ORTHODOX]}: minimal coupling is the standard gateway from free to gauge-coupled wave dynamics.

\textbf{[SPECULATIVE]}: the above substitution is only an analogy template unless a full gauge-covariant derivation is supplied.

\subsection{Ribbon-Invariant Chirality Refinement}
The older canon layers also suggested using the Calugareanu--White--Fuller relation
\begin{align}
    Lk = Tw + Wr
\end{align}
as a refinement of the knot-based matter/antimatter dictionary.

\textbf{[ORTHODOX]}: this is a standard geometric identity for framed curves and ribbons.

\textbf{[DERIVED]}: it can be used internally to separate writhe-dominated from twist-dominated realizations of the same coarse knot class.

\textbf{[SPECULATIVE]}: the mapping from $\operatorname{sign}(Tw)$ or related chirality measures to particle/antiparticle sectors remains a model hypothesis.

\subsection{Particle-Candidate Mapping Layer}
\textbf{[SPECULATIVE] Working dictionary:} a minimal knot/composition mapping retained from earlier SST work is summarized in Table~\ref{tab:particle-candidates}. This table is organizational rather than definitive; it records the current topological assignments used by the mass and taxonomy program.

\begin{table}[h]
\centering
\begin{tabular}{lll}
\hline
Knot class / composition & Candidate & Status / role \\
\hline
$3_1$ & electron & baseline lepton calibration \\
$T(2,5)$ or equivalent higher torus class & muon & lepton hierarchy candidate \\
$T(2,7)$ or equivalent higher torus class & tau & lepton hierarchy candidate \\
$5_2 + 5_2 + 6_1$ & proton & composite baryon candidate \\
analogous nearby composite class & neutron & composite baryon candidate \\
\hline
\end{tabular}
\caption{Minimal knot-class to particle-candidate summary retained as a research-track dictionary.}
\label{tab:particle-candidates}
\end{table}

\textbf{[CRITICAL NOTE]}: Table~\ref{tab:particle-candidates} is a working dictionary only. A full canonical assignment requires an explicit invariant set and a mass-functional comparison program.

\subsection{Hydrodynamic Exchange and the Pauli Barrier}
\textbf{[ORTHODOX]}: for two thin filamentary loops $C_1$ and $C_2$ in an incompressible medium, the Biot--Savart interaction energy is
\begin{align}
    E_{\mathrm{int}}
    =
    \frac{\rhoF\Gamma_1\Gamma_2}{4\pi}
    \oint_{C_1}\oint_{C_2}
    \frac{d\ell_1\cdot d\ell_2}{\norm{r_1-r_2}}.
\end{align}

\textbf{[DERIVED]}: if two identical electron-candidate loops are forced into the same spatial state, regularization at the core radius $\rc$ yields the overlap penalty
\begin{align}
    V_{\mathrm{Pauli}}(d)
    \approx
    \frac{\rhoF\Gamma_0^2}{4\pi}
    \iint
    \frac{d\ell_1\cdot d\ell_2}{\sqrt{\norm{r_1(\theta_1)-r_2(\theta_2)+d}^2+\rc^2}},
\end{align}
with maximal-overlap scale
\begin{align}
    V_{\mathrm{Pauli}}^{\max}
    \approx
    \frac{\rhoF\Gamma_0^2}{4\pi}\left(\frac{L}{\rc}\right)\mathcal{O}(1).
\end{align}

Using the canonical values
\begin{align}
    \rhoF &= 7.0\times 10^{-7}\ \mathrm{kg\,m^{-3}}, \\
    \rc &= 1.40897017\times 10^{-15}\ \mathrm{m}, \\
    \Gamma_0 &= 2\pi \rc\,\vnorm \approx 9.6836\times 10^{-9}\ \mathrm{m^2\,s^{-1}},
\end{align}
and $L\sim 2\pi a_0$ at the hydrogenic scale, one obtains the benchmark estimate
\begin{align}
    V_{\mathrm{Pauli}}^{\max} \sim 1.23\times 10^{-18}\ \mathrm{J} \sim 7.6\ \mathrm{eV}.
\end{align}

\textbf{[CRITICAL NOTE]}: this is retained as a canonical scale benchmark, not yet as a closed theorem. Its role is to show that the exclusion-energy estimate falls in an atomically relevant window rather than at obviously pathological scales.

\subsection{Numerical Benchmark Layer}
\textbf{[DERIVED] Computational-program requirement:} earlier SST work contained explicit benchmark tables for particle and atomic masses. In v0.8.1, the benchmark layer is retained as a reproducibility requirement rather than as a stand-alone authority claim.

\begin{table}[h]
\centering
\begin{tabular}{lllll}
\hline
Particle & $m_{\mathrm{exp}}$ (MeV) & $m_{\mathrm{SST}}$ (MeV) & rel. error & mode class \\
\hline
e & 0.51099895 & 0.51099895 & 0 & predictive/closure \\
$\mu$ & 105.6583745 & \dots & \dots & predictive candidate \\
$\tau$ & 1776.86 & \dots & \dots & predictive candidate \\
p & 938.272088 & \dots & \dots & predictive or closure \\
n & 939.565420 & \dots & \dots & predictive or closure \\
\hline
\end{tabular}
\caption{Benchmark summary structure to be populated from canonical scripts and archived CSV output.}
\end{table}

\textbf{[CRITICAL NOTE]}: a future v0.8.x benchmark appendix should include the exact script or package path used for generation, archived CSV outputs, mode definitions (predictive vs. exact-closure), and uncertainty propagation on canonical constants and geometric factors.

\subsection{Helicity-by-base archive (diagnostic layer)}
\textbf{[DERIVED] Archive role:} the SST helicity-by-base outputs are retained as a measurable balance diagnostic for taxonomy support, not as an independent law.
Define the midpoint-distance diagnostic
\begin{align}
    \Delta_H(K)=\abs{\widehat{\mathcal{H}}_{b}(K)+\frac{1}{2}},
\end{align}
where $\widehat{\mathcal{H}}_{b}(K)$ is the archived normalized helicity-like base indicator for class $K$.

Near-balanced values with $\Delta_H(K)\ll 1$ indicate candidate dark or quasi-dark sectors only when symmetry and instability filters are also satisfied.

\textbf{[CRITICAL NOTE]}: no single helicity-like scalar is sufficient for sector assignment; classifier decisions remain conjunction rules over symmetry, helicity balance, and low-order stability content.

\textbf{[CANON WORKFLOW NOTE]}: if SST-labelled and parallel VAM-labelled helicity outputs are numerically identical, the SST-labelled archive is the only required source for canon workflows.

\subsection{Trefoil-closure archive (benchmark discipline)}
\textbf{[DERIVED] Benchmark role:} ideal-trefoil robustness sweeps, final-best-estimate tables, and theorem-style reports are retained as benchmark-discipline infrastructure.

\textbf{[CRITICAL NOTE]}: practical best-estimate closure and theorem-level closure are distinct.
A result may remain useful for calibration while still failing strict theorem-level tail or postcheck requirements.

\textbf{[REQUIRED METADATA]}: benchmark reports should include explicit pass/fail logic, postchecks, extrapolation summaries, and local-curve diagnostics so that closure claims are auditable.

\subsection{\texorpdfstring{$\phi_{3_1}$}{phi3_1} transform archive (deferred)}
\textbf{[SPECULATIVE] Status:} retained as research-track spectral evidence for later zeta/spectral programs.

\textbf{[SCOPE RULE]}: this archive is not part of the current dark-sector or galactic canonization closure path in v0.8.x.

\subsection{Gauge-Sector Roadmap}
\textbf{[SPECULATIVE] Structural roadmap:} a retained minimal statement of the gauge program is
\begin{align}
    g_{\mathrm{eff}} = su(3) \oplus su(2) \oplus u(1),
\end{align}
with effective gauge phases organized as holonomy data of multi-director transport.

A minimal phase-based representation is
\begin{align}
    A^a_{\mu}(x) = \partial_{\mu}\theta^a(x),
\end{align}
and clock coupling acts as a common phase reference,
\begin{align}
    \theta^a(x) \mapsto \theta^a(x)+q^a\chi(x)
    \quad\Rightarrow\quad
    A^a_{\mu}\mapsto A^a_{\mu}+q^a\partial_{\mu}\chi.
\end{align}

\textbf{[CRITICAL NOTE]}: the gauge layer is retained only as a roadmap statement. Running couplings, anomaly cancellation, and realistic mixing matrices remain open work.

\section{Dark-sector and galactic canonization execution package (v0.8.x)}
\label{sec:dark-sector-execution-package}

\subsection{Stage-0 freeze: topology toolkit and benchmark protocol}
\textbf{[DERIVED] Scope freeze:} the mandatory topology toolkit for v0.8.x canonization is fixed to:
\begin{enumerate}
    \item Gauss linking integral,
    \item helicity in ideal fluids,
    \item thin-tube helicity reduction,
    \item Hopf invariant,
    \item basic torus-knot invariants $(c,b,g)$,
    \item Calugareanu--White--Fuller relation $Lk=Tw+Wr$.
\end{enumerate}

\textbf{[DERIVED] Notation freeze:} one notation family is used in this package:
\[
\rhoF,\quad v_\theta(r),\quad p_{\mathrm{swirl}}(r),\quad
\mathcal{H}(K),\quad N_u^{(1)}(K),\quad \Gamma,\quad \kappa,\quad \SwirlClock.
\]
No parallel symbol aliases are introduced for the same quantity.

\textbf{[DERIVED] Mandatory benchmark protocol template:}
\begin{enumerate}
    \item \textbf{Inputs:} canonical constants, topology identifiers, profile parameters, and dataset/version tags.
    \item \textbf{Procedure:} exact script or package path, run command, and branch/commit hash.
    \item \textbf{Outputs:} archived CSV files with fixed column schema and an attached metadata note.
    \item \textbf{Uncertainty:} first-order propagation for calibrated constants and sensitivity sweep on profile parameters.
    \item \textbf{Rerun:} deterministic rerun steps and acceptance tolerances.
\end{enumerate}

\textbf{[CANON WORKFLOW NOTE]}: every promoted quantitative claim in this package is invalid without an explicit input-to-CSV trace.

\subsection{Stage-1 freeze: dark-sector Euler pressure anchor}
\textbf{[ORTHODOX] Euler radial balance (steady, axisymmetric, azimuthal-dominant):}
\begin{align}
    \frac{1}{\rhoF}\frac{dp_{\mathrm{swirl}}}{dr}
    =
    \frac{v_\theta(r)^2}{r}.
    \label{eq:exec_swirl_pressure}
\end{align}

\textbf{[DERIVED] Flat-tail limit (constant tail speed):}
\begin{align}
    v_\theta(r)\to v_0
    \quad\Longrightarrow\quad
    p_{\mathrm{swirl}}(r)
    =
    p_0 + \rhoF v_0^2\ln\!\left(\frac{r}{r_0}\right).
\end{align}

\textbf{[DERIVED] Assumption block:}
incompressible, inviscid, stationary, axisymmetric, azimuthal-dominant flow, with validity restricted to regions where radial drift and stratification corrections are subleading.

\textbf{[DERIVED] Dimensional check:}
\[
\left[\frac{1}{\rhoF}\frac{dp}{dr}\right]
=
\frac{\mathrm{kg\,m^{-2}\,s^{-2}}}{\mathrm{kg\,m^{-3}}}
=
\mathrm{m\,s^{-2}}
=
\left[\frac{v^2}{r}\right].
\]

\textbf{[SPECULATIVE] SST interpretation:}
Eq.~\eqref{eq:exec_swirl_pressure} is interpreted as a classifier/driver layer for dark-sector support. This interpretation is model-level and separate from the orthodox derivation.

\textbf{[DERIVED] Attached observables/falsifiers:}
\begin{enumerate}
    \item Rotation-curve consistency after pressure reconstruction from archived $v_\theta(r)$ data.
    \item Null falsifier: systematic mismatch in the reconstructed pressure gradient sign or scale across validated radial windows.
\end{enumerate}

\subsection{Stage-2 freeze: measurable dark taxonomy}
\textbf{[DERIVED] Minimal classifier variables:}
symmetry class (including amphichirality status), chirality sign/state, helicity-balance indicator $\Delta_H(K)$, and low-order instability count $N_u^{(1)}(K)$.

\textbf{[DERIVED] Decision rules (symmetry-first, then stability-filtered):}
\begin{align}
    \text{Dark}
    &:\ \text{amphichiral candidate}
    \ \wedge\
    \Delta_H(K)\ll 1
    \ \wedge\
    N_u^{(1)}(K)=0,\\
    \text{Quasi-dark}
    &:\ \text{near-balanced symmetry/helicity}
    \ \wedge\
    0 < N_u^{(1)}(K)\le 2,\\
    \text{Visible}
    &:\ \text{chiral-dominant and/or helicity-unbalanced sectors}.
\end{align}

\textbf{[DERIVED] Ambiguity tie-break:}
if classes overlap, choose the class with the strongest conjunction score in the order
\[
\text{symmetry gate} \rightarrow \text{helicity balance gate} \rightarrow \text{stability gate}.
\]

\textbf{[DERIVED] Baseline templates and falsifiers:}
\begin{itemize}
    \item \textbf{Visible baseline:} trefoil-like chiral template; falsifier = persistent amphichiral-plus-balanced signature under validated data updates.
    \item \textbf{Quasi-dark baseline:} Borromean-type weak-coupling template; falsifier = stable zero-instability assignment across independent scans.
    \item \textbf{Dark baseline:} figure-eight/amphichiral template; falsifier = robust nonzero chirality or high-instability assignment.
\end{itemize}

\textbf{[CANON WORKFLOW NOTE]}:
amphichiral $\Rightarrow$ dark remains a classifier proposal tested against observables, not a closed theorem.

\subsection{Stage-3 freeze: galactic law branch gate and closure route}
\textbf{[DERIVED] Branch-gate decision:} Branch A (legacy profile refine) is selected for current v0.8.x closure because it preserves tested limits while allowing benchmark traceability.

Working profile:
\begin{align}
    v_\theta(r)
    =
    \frac{C_{\mathrm{core}}}{\sqrt{1+(r_c/r)^2}}
    +
    C_{\mathrm{tail}}\left(1-e^{-r/r_c}\right).
\end{align}

\textbf{[DERIVED] Required checks:}
small-$r$ behavior, flat-tail behavior, parameter semantics, and pressure reconstruction via Eq.~\eqref{eq:exec_swirl_pressure}.

\textbf{[DERIVED] Failure modes:}
\begin{enumerate}
    \item non-physical parameter combinations (wrong sign or unstable growth),
    \item failure to satisfy tail-limit behavior in benchmark windows,
    \item inability to reconstruct a consistent pressure profile from the same archived inputs.
\end{enumerate}

\subsection{Stage-4 freeze: probe falsification map}
\textbf{[DERIVED]} probe classes are fixed as astrophysical, laboratory pressure/swirl analog, HV-coil/thrust-inspired, optical/torsional coupling, and null dark-sector tests.

\begin{table}[h]
\centering
\begin{tabular}{p{2.7cm}p{3.1cm}p{2.5cm}p{2.2cm}p{2.2cm}}
\hline
Probe class & Equation under test & Null-result meaning & Primary confounder & Expected sign/trend \\
\hline
Astrophysical rotation curves & Eq.~\eqref{eq:exec_swirl_pressure} + profile closure & weakens pressure-law applicability window & baryonic-model bias & reconstructed $dp/dr$ consistent with $v_\theta^2/r$ \\
Lab swirl analog & Euler radial balance scaling & weakens scale-transfer claim & viscosity/edge forcing & monotone pressure rise with imposed swirl \\
HV-coil/thrust-inspired & circulation-pressure bridge claims & kills quantized-bridge claim if no scaling & EM/thermal artifacts & response tracks circulation proxy \\
Optical/torsional coupling & swirl-clock bridge signatures & weakens coupling interpretation & detector drift/noise & phase/timing shift with controlled swirl proxy \\
Null dark tests & taxonomy conjunction rules & kills classifier branch if conjunction fails & mislabelled topology class & class score decreases under contradictory diagnostics \\
\hline
\end{tabular}
\caption{Frozen equation-to-probe map for falsification-first workflow in the v0.8.x dark-sector program.}
\label{tab:exec-probe-map}
\end{table}

\textbf{[CANON WORKFLOW NOTE]}: probe outcomes can support or reject already-stated equations, but they do not replace derivation-level closure.

\subsection{Promotion snapshot after execution}
\textbf{[DERIVED]} post-freeze status in this package:
\begin{itemize}
    \item Item 6 (topology toolkit): CANON-ready in scope and notation.
    \item Item 5 (benchmark protocol): CANON-CANDIDATE pending repository-wide protocol reuse.
    \item Item 1 (Euler pressure anchor): CANON-CANDIDATE with explicit assumptions, limits, and falsifier path.
    \item Item 2 (dark taxonomy): BRIDGE with measurable classifier gates and falsifiers.
    \item Item 3 (galactic swirl law): BRIDGE with branch gate closed and benchmark route fixed.
    \item Item 4 (probe map): CANON-BRIDGE with equation/null-result mapping frozen.
\end{itemize}
